\section{The Schmidt number and the full reduction criterion}

The Schmidt number of a bipartite state is the smallest $n$ for which the state has
a convex decomposition into pure states each of Schmidt rank at most $n$. Separable
states are precisely those of Schmidt number one, and a state of Schmidt number at
most $n$ satisfies the reduction criterion
\begin{align}
    n\,\rho_1\otimes\Id & \ge\rho,
    \notag                         \\
    n\,\Id\otimes\rho_2 & \ge\rho,
    \notag
\end{align}
which is \cite[Equation~(3.18)]{Wolf2012Quantum}. The bound on the Schmidt number is
recorded as a finite sum of pure-state projectors of Schmidt rank at most $n$, the
unnormalized form of the convex decomposition; rescaling a vector by the square root
of its weight scales its coefficient matrix and so leaves the Schmidt rank unchanged.

\begin{definition}[Bounded Schmidt number]
    \label{def:has_schmidt_number_le}
    \lean{Matrix.HasSchmidtNumberLE}
    \leanok
    \uses{def:schmidt_rank}
    A bipartite matrix $\rho$ on $\MN{d}\otimes\MN{d'}$ has Schmidt number at most $n$
    when it is a finite sum of pure-state projectors of Schmidt rank at most $n$,
    \begin{align}
        \rho        & =\sum_i\ket{\psi_i}\bra{\psi_i},
        \notag                                         \\
        \SR(\psi_i) & \le n.
        \notag
    \end{align}
\end{definition}

The closure and compactness properties of the Schmidt-number sets are proved
in Section~\ref{sec:app_positive_not_cp_schmidt_number_geometry}.

\begin{theorem}[Schmidt number one is separability]
    \label{thm:has_schmidt_number_le_one_iff_separable}
    \lean{Matrix.hasSchmidtNumberLE_one_iff_isSeparable}
    \leanok
    \uses{def:has_schmidt_number_le, def:is_separable, def:schmidt_rank,
        lem:matrix_rank_factorization_coordinate, thm:schmidt_rank_basic_bounds,
        thm:is_separable_kronecker, thm:is_separable_add}
    A bipartite matrix has Schmidt number at most one if and only if it is separable.
\end{theorem}

\begin{proof}
    \leanok
    A pure state of Schmidt rank at most one has a coefficient matrix of rank at most
    one, which factors as $\psi(i,j)=u_i v_j$; its projector is the Kronecker product
    $\ket{u}\bra{u}\otimes\ket{v}\bra{v}$ of two rank-one positive semidefinite
    matrices, hence separable, and separable matrices are closed under addition.
    Conversely a Kronecker product $A\otimes B$ of positive semidefinite matrices,
    after expanding each factor into rank-one projectors, is a double sum of
    product-vector projectors, each of Schmidt rank at most one, so it has Schmidt
    number at most one; a separable matrix is a finite sum of such products.
\end{proof}

\begin{theorem}[Pure-state step of the reduction criterion]
    \label{thm:tensor_map_id_t_eta_posSemidef_pure}
    \lean{Matrix.tensorMapId_tEta_posSemidef_of_hasSchmidtRankLE}
    \leanok
    \uses{def:schmidt_rank, def:t_eta, def:tensor_map_id, thm:t_eta_threshold,
        thm:square_schmidt_rank_exact_parametrization,
        lem:rank_one_ampliation_choi_compression,
        thm:choi_right_compression_quadratic_form, thm:schmidt_rank_choi_test_iff}
    For a pure state $\ket{\psi}\bra{\psi}$ on the square bipartite system
    $\MN{D}\otimes\MN{D}$ with $\SR(\psi)\le n$ and
    $1\le n<D$, applying $T_n$ to the first factor keeps the result positive
    semidefinite: $(T_n\otimes\id)(\ket{\psi}\bra{\psi})\ge 0$.
\end{theorem}

\begin{proof}
    \leanok
    The map $T_n$ is $n$-positive for $1\le n<D$. Write $\psi$ through the maximally
    entangled vector as $\psi=\psi_X$ for a square matrix $X$ on the right factor with
    $\rank(X)=\SR(\psi)\le n$. Then
    $(T_n\otimes\id)(\ket{\psi}\bra{\psi})$ is the right-factor Choi compression of
    $T_n$ by $X$. Its quadratic form on any vector $\eta$ is the Choi quadratic form
    of $T_n$ on $R_X^\dagger\eta$, whose Schmidt rank is at most
    $\rank(X)\le n$. The Schmidt-rank Choi criterion makes that form
    non-negative, so the compression is positive semidefinite.
\end{proof}

\begin{theorem}[Pure-state step for a general positive map]
    \label{thm:tensor_map_id_posSemidef_pure}
    \lean{Matrix.tensorMapId_posSemidef_of_hasSchmidtRankLE}
    \leanok
    \uses{def:schmidt_rank, def:tensor_map_id, def:k_positive_map,
        def:choi_right_compression, thm:square_schmidt_rank_exact_parametrization,
        lem:rank_one_ampliation_choi_compression,
        thm:choi_right_compression_quadratic_form, thm:schmidt_rank_choi_test_iff}
    For a pure state $\ket{\psi}\bra{\psi}$ on the square bipartite system
    $\MN{D}\otimes\MN{D}$ with $\SR(\psi)\le n$ and any $n$-positive map
    $T$, applying $T$ to the first factor keeps the result positive semidefinite:
    $(T\otimes\id)(\ket{\psi}\bra{\psi})\ge 0$.
\end{theorem}

\begin{proof}
    \leanok
    Write $\psi$ through the maximally entangled vector as $\psi=\psi_X$ for a square
    matrix $X$ on the right factor with
    $\rank(X)=\SR(\psi)\le n$. Then
    $(T\otimes\id)(\ket{\psi}\bra{\psi})$ is the right-factor Choi compression of $T$
    by $X$, and its quadratic form on a vector $\eta$ equals the Choi quadratic form
    $C_T$ of $T$ evaluated on $R_X^\dagger\eta$, a vector of Schmidt rank at most
    $\rank(X)\le n$. Because $T$ is $n$-positive,
    $\langle R_X^\dagger\eta,\, C_T\, R_X^\dagger\eta\rangle\ge 0$, so the
    compression is positive semidefinite.
\end{proof}

\begin{theorem}[Positive maps and entanglement, only-if direction]
    \label{thm:has_schmidt_number_le_tensor_map_id}
    \lean{Matrix.HasSchmidtNumberLE.tensorMapId_posSemidef}
    \leanok
    \uses{def:has_schmidt_number_le, def:tensor_map_id, def:k_positive_map,
        thm:tensor_map_id_sum, thm:tensor_map_id_posSemidef_pure}
    A bipartite state on the square system $\MN{D}\otimes\MN{D}$ of Schmidt number at
    most $n$ satisfies $(T\otimes\id)(\rho)\ge 0$ for every $n$-positive map $T$.
    This is the square case $d=d'=D$ of the only-if direction of Wolf's
    Proposition~3.4 for a general bipartite system
    $\MN{d}\otimes\MN{d'}$ \cite[Chapter~3, Proposition~3.4]{Wolf2012Quantum}.
\end{theorem}

\begin{proof}
    \leanok
    The state is a finite sum of pure-state projectors of Schmidt rank at most $n$.
    Applying $T$ to the first factor is linear, so it distributes over the sum; the
    pure-state step makes each summand positive semidefinite, and a finite sum of
    positive semidefinite matrices is positive semidefinite.
\end{proof}

\begin{theorem}[Only-if direction for a general bipartite system]
    \label{thm:has_schmidt_number_le_tensor_map_id_general}
    \lean{Matrix.HasSchmidtNumberLE.tensorMapId_posSemidef_general}
    \leanok
    \uses{def:has_schmidt_number_le, def:tensor_map_id, def:k_positive_map,
        thm:has_schmidt_number_le_tensor_map_id}
    On a general bipartite system $\MN{d}\otimes\MN{d'}$, with no relation imposed
    between the two factors, a state of Schmidt number at most $n$ satisfies
    $(T\otimes\id)(\rho)\ge 0$ for every $n$-positive map $T$ on $\MN{d}$.
    This is the only-if direction of Wolf's Proposition~3.4, the first step of
    eq.~(3.18) \cite[Chapter~3, Proposition~3.4]{Wolf2012Quantum}, with no restriction
    on the bipartite dimensions.
\end{theorem}

\begin{proof}
    \leanok
    Pad the bipartite state to the square system $\MN{D}\otimes\MN{D}$ with
    $D=\max(d,d')$ by zero rows and columns: when $d'\le d$ the second factor is padded
    up to $d$; when $d\le d'$ the first factor is padded up to $d'$ and $T$ is extended
    to a map on $\MN{d'}$ that preserves $n$-positivity. Padding preserves the
    Schmidt-number bound, the square result makes the padded ampliation positive
    semidefinite, and the original ampliation is its principal submatrix, hence positive
    semidefinite.
\end{proof}

\begin{theorem}[Reduction step from the Schmidt-number premise]
    \label{thm:has_schmidt_number_le_tensor_map_id_t_eta}
    \lean{Matrix.HasSchmidtNumberLE.tensorMapId_tEta_posSemidef}
    \leanok
    \uses{def:has_schmidt_number_le, def:t_eta, def:tensor_map_id,
        thm:tensor_map_id_sum, thm:tensor_map_id_t_eta_posSemidef_pure}
    For $1\le n<D$, a bipartite state on the square system
    $\MN{D}\otimes\MN{D}$ of Schmidt number at most $n$ satisfies
    $(T_n\otimes\id)(\rho)\ge 0$.
\end{theorem}

\begin{proof}
    \leanok
    The state is a finite sum of pure-state projectors of Schmidt rank at most $n$.
    Applying $T_n$ to the first factor is linear, so it distributes over the sum; the
    pure-state step makes each summand positive semidefinite, and a finite sum of
    positive semidefinite matrices is positive semidefinite.
\end{proof}

\begin{theorem}[Reduction criterion with the Schmidt-number premise]
    \label{thm:reduction_criterion_of_schmidt_number}
    \lean{Matrix.reductionCriterion_left_of_hasSchmidtNumberLE_general,
        Matrix.reductionCriterion_right_of_hasSchmidtNumberLE_general}
    \leanok
    \uses{def:has_schmidt_number_le, def:t_eta,
        thm:has_schmidt_number_le_tensor_map_id_general, thm:reduction_criterion}
    On a general bipartite system $\MN{d}\otimes\MN{d'}$, a state of Schmidt number at
    most $n$ satisfies the reduction criterion
    \begin{align}
        n\,\Id\otimes\rho_2 & \ge\rho
                            &         & (1\le n<d),
        \notag                                       \\
        n\,\rho_1\otimes\Id & \ge\rho
                            &         & (1\le n<d'),
        \notag
    \end{align}
    with $\rho_1$ and $\rho_2$ the reduced densities on the first and second factors.
    No relation is imposed between $d$ and $d'$; each inequality carries only the
    $n$-positivity threshold of $T_n$ on its acting factor.
\end{theorem}

\begin{proof}
    \leanok
    The Schmidt-number premise gives $(T_n\otimes\id)(\rho)\ge 0$ on the general system
    (first step), and the dimension-general operator implication then yields
    $n\,\Id\otimes\rho_2\ge\rho$. The factor swap of a state of Schmidt number at most
    $n$ again has Schmidt number at most $n$, since it sends each pure summand $\psi$ to
    $\psi\circ\mathrm{swap}$, whose Schmidt coefficient matrix is the transpose
    $M_{\psi\circ\mathrm{swap}}=M_\psi^{\mathsf T}$ of that of $\psi$ and so has the same
    rank; applying the first step to the swapped state (whose first factor is $d'$) and
    the operator implication in its symmetric form gives $n\,\rho_1\otimes\Id\ge\rho$.
\end{proof}


The separating-hyperplane construction and its trace-form representation are
proved in Section~\ref{sec:app_positive_not_cp_schmidt_number_geometry}.

\begin{theorem}[Witness criterion for Schmidt number]
    \label{thm:not_has_schmidt_number_le_iff_exists_witness}
    \lean{Matrix.not_hasSchmidtNumberLE_iff_exists_witness}
    \leanok
    \uses{def:has_schmidt_number_le, def:schmidt_rank,
        thm:exists_entanglement_witness}
    Let $\rho$ be a trace-one Hermitian bipartite state. Then the Schmidt number of
    $\rho$ exceeds $n$ if and only if there is a Hermitian operator $W$ such that, for
    every $\psi$ of Schmidt rank at most $n$,
    \begin{align}
        \re\tr(W\rho)            & <0,
        \notag                            \\
        \re\bra{\psi}W\ket{\psi} & \ge 0.
        \notag
    \end{align}
    This is Wolf's Proposition~3.3 \cite[Chapter~3, Proposition~3.3]{Wolf2012Quantum}: a
    state has Schmidt number larger than $n$ exactly when it is detected by an
    entanglement witness for $S_n$.
\end{theorem}

\begin{proof}
    \leanok
    The forward implication is the separating-hyperplane construction
    (\ref{thm:exists_entanglement_witness}). For the converse, suppose $\rho$ had Schmidt
    number at most $n$, so $\rho=\sum_i\ket{\psi_i}\bra{\psi_i}$ with each $\psi_i$ of
    Schmidt rank at most $n$. By linearity of the trace,
    $\re\tr(W\rho)=\sum_i\re\tr\!(W\ket{\psi_i}\bra{\psi_i})$,
    a sum of non-negative terms, contradicting $\re\tr(W\rho)<0$. This
    converse uses no density-matrix hypotheses on $\rho$.
\end{proof}

The Choi correspondence and trace-pairing identities used below are proved in
Section~\ref{sec:app_positive_not_cp_witness_choi}.

\begin{theorem}[Positive maps detect high Schmidt number]
    \label{thm:exists_isNPositiveMap_detects}
    \lean{Matrix.exists_isNPositiveMap_tensorMapId_not_posSemidef}
    \leanok
    \uses{def:k_positive_map, def:has_schmidt_number_le,
        thm:exists_entanglement_witness, thm:exists_isNPositiveMap_choiMatrix_eq,
        thm:k_positive_trace_pairing_adjoint, thm:trace_choiMatrix_omega_quadratic_adjoint}
    Let $\rho$ be a trace-one Hermitian bipartite state on $\C^D\otimes\C^D$ whose Schmidt
    number exceeds $n$. Then there is an $n$-positive map $T\colon\MN{D}\to\MN{D}$ such that
    $(T\otimes\Id)(\rho)$ is not positive semidefinite. This is the if direction of Wolf's
    Proposition~3.4 \cite[Chapter~3, Proposition~3.4]{Wolf2012Quantum}.
\end{theorem}

\begin{proof}
    \leanok
    The entanglement witness $W$ for $\rho$ (\ref{thm:exists_entanglement_witness}) is the
    Choi matrix of an $n$-positive map $P$ (\ref{thm:exists_isNPositiveMap_choiMatrix_eq}).
    Its trace-pairing adjoint $T=P^{*}$ is again $n$-positive
    (\ref{thm:k_positive_trace_pairing_adjoint}). By the Choi trace pairing through the
    trace-pairing adjoint (\ref{thm:trace_choiMatrix_omega_quadratic_adjoint}),
    $\bra{\Omega}(T\otimes\Id)(\rho)\ket{\Omega}=\tr(W\rho)$, whose real part is
    negative. A positive semidefinite matrix has non-negative quadratic forms, so
    $(T\otimes\Id)(\rho)$ is not positive semidefinite.
\end{proof}
